\documentclass[11pt]{tex/lawoutline}
\usepackage{amssymb}
\usepackage{enumitem}

\course{Legislation and Regulation}
\student{Prof: R. Doerfler}
\semester{Fall 2026}

\begin{document}

\maketitle
\tableofcontents

\chapter{In-Class Problems}

\section{Sep 23}

\lincanon{Rule of Last Antecedent}
{A limiting phrase ordinarily modifies only the noun or phrase immediately preceding it.}

\lincanon{Series-Qualifier Principle}
{When a modifier follows a series of parallel terms, the modifier ordinarily applies to every term in the series when that reading is natural.}

\subcanon{Rule of Lenity}
{Ambiguity in a criminal statute should generally be resolved in favor of the defendant.}

\casebox{\textit{Lockhart v. United States}}
{Supreme Court, 2016}
{
\textbf{Facts}
\begin{itemize}
    \item Lockhart was convicted of sexual abuse in the 1st degree, which involved his adult girlfriend. 
    \item Separately, Lockhart was indicted (and pleaded guilty) for attempting to receive child pornography that violated 18 U.S.C Section 2252(a)(2) and for possession of child pornography, which violated Section 2252(a)(4)(b). 
    \item Presentence report concluded that Lockhart was subject to Section 2252(b)(2)'s mandatory minimum because of his prior conviction related to aggravated sexual abuse, sexual abuse, or abusive sexual conduct involving a minor or ward. 
    \item Here, basically saying that minor/ward part of the statute only applies to the mention of abusive sexual conduct right before, which would construe the statute quite generally.
    \item Lockhart argued that minor/ward applied to all three of the listed crimes, and that his charges of sexual abuse involving an adult fell outside the scope. 
\end{itemize}

\rulebox{Section 2252(b)(2)}
{Defendants convicted of possessing child pornography in violation of 18 U.S.C. § 2252(a)(4) are subject to a 10–year mandatory minimum sentence and an increased maximum sentence if they have ``a prior conviction ... under the laws of any State relating to aggravated sexual abuse, sexual abuse, or abusive sexual conduct involving a minor or ward.''} 

\textbf{Question} Does the phrase ``involving a minor or ward'' modify all items in the list of predicate crimes, or only the item that precede it?

\textbf{Holding}
\begin{itemize}
    \item Held that the minor/ward clause only modifies the ``abusive sexual conduct'' phrase directly in front of it. 
    \item Therefore, Lockhart's previous indictment does apply in this case. 
\end{itemize}

\textbf{Discussion}

\begin{enumerate}
    \item \textbf{Rule of last antecedent vs Series-qualifier principle}
    \begin{itemize}
        \item \textbf{Majority:} the modifier applies only to ``abusive sexual conduct.''
        \item \textbf{Dissent:} the statute contains an integrated list, so the series-qualifier principle applies.
    \end{itemize}

    \item \textbf{Whether context overcomes the grammatical presumption}
    \begin{itemize}
        \item \textbf{Majority:} the list is structurally irregular, and Chapter 109A supports the last-antecedent reading.
        \item \textbf{Dissent:} the offenses are sufficiently parallel and related that ordinary English applies the modifier to all three.
    \end{itemize}

    \item \textbf{Statutory structure and Chapter 109A}
    \begin{itemize}
        \item \textbf{Majority:} Congress modeled the state predicates on Chapter 109A, which includes offenses involving adults and minors.
        \item \textbf{Majority:} Wouldn't make sense for there to be two different outcomes of almost identical state and federal crimes. 
        \item \textbf{Dissent:} The statutes do not precisely match; Chapter 109A contains four offenses, while Section 2252(b)(2) lists only three.
        \item \textbf{Dissent:} The majority placed too much weight on the similarities between the statutes, but the categories do not precisely mirror one another, and so cannot say Congress intended to reproduce Chapter 109A's structure in Section 2252(b)(2).
    \end{itemize}

    \item \textbf{Rule against surplusage}
    \begin{itemize}
        \item \textbf{Majority:} applying the modifier to all three offenses makes the list excessively redundant.
        \item \textbf{Dissent:} the majority's interpretation also creates redundancy and makes ``involving a minor or ward'' largely unnecessary.
    \end{itemize}

    \item \textbf{Rule of lenity}
    \begin{itemize}
        \item \textbf{Majority:} lenity applies only after ordinary interpretive tools fail; here, the text and structure provide a satisfactory interpretation.
        \item \textbf{Dissent:} because the statute remains reasonably ambiguous and imposes a criminal mandatory minimum, lenity should favor Lockhart.
    \end{itemize}
\end{enumerate}
}

\end{document}